% BHU PYQ -- Professional Exam Template
\documentclass[11pt,a4paper]{article}

\usepackage[a4paper,top=2.5cm,bottom=2.5cm,left=2.8cm,right=2.8cm,headheight=14pt]{geometry}
\usepackage{amsmath,amssymb}
\usepackage[T1]{fontenc}
\usepackage[utf8]{inputenc}
\usepackage{lmodern}
\usepackage{microtype}
\usepackage{fancyhdr}
\usepackage{enumitem}
\usepackage{hyperref}
\usepackage{lastpage}
\usepackage{parskip}

\hypersetup{colorlinks=true,urlcolor=black,linkcolor=black,
  pdftitle={B.Sc. (Hons.) Botany Sem-II BOB-201 2013-14 -- BHU PYQ}}

\pagestyle{fancy}
\fancyhf{}
\renewcommand{\headrulewidth}{0.4pt}
\renewcommand{\footrulewidth}{0.4pt}
\fancyhead[L]{\small   \textit{Banaras Hindu University}}
\fancyhead[R]{\small   \textit{Botany Paper: BOB-201 (2013-14)}}
\fancyfoot[C]{\small Page  \thepage\ of \pageref{LastPage}}


\newlist{parts}{enumerate}{2}
\setlist[parts,1]{label=(\alph*),leftmargin=*,itemsep=4pt,topsep=3pt}
\setlist[parts,2]{label=(\roman*),leftmargin=*,itemsep=2pt,topsep=2pt}

\newcommand{\pts}[1]{\hfill   \textit{[#1]}}


\begin{document}


\begin{center}
  {\large\bfseries BANARAS HINDU UNIVERSITY}\\[2pt]
  {
\normalsize B.Sc. (Hons.) II Semester Examination, 2013-14}\\[6pt]
  \rule{0.8\linewidth}{0.5pt}\\[6pt]
  {\large\bfseries Botany}\\[4pt]
  {\large\bfseries Paper: BOB-201 --- Microbiology, Plant Pathology, Cytology and Genetics}\\[4pt]
  {
\normalsize Time: Three Hours \hfill Maximum Marks: 70}
\end{center}

\rule{\linewidth}{0.4pt}

\smallskip



\noindent   \textit{Note: Attempt five questions in all, selecting two questions each from Section A and B. Question No. 9 (Section C) is compulsory. Marks allotted to each question are given on the right.}

\bigskip


\begin{center}
   \textbf{SECTION -- A}
\end{center}
\medskip

\medskip


\noindent   \textbf{Question 1.} With the help of suitable diagrams describe the structure of the bacterial cell wall and flagella.\pts{15}

\medskip


\noindent   \textbf{Question 2.} Write short notes on any two of the following:\pts{$7.5  \times 2 = 15$}

\begin{parts}
  \item Plasmids
  \item Endospore
  \item Lytic cycle
\end{parts}

\medskip


\noindent   \textbf{Question 3.} Write critical notes on any two of the following:\pts{$7.5  \times 2 = 15$}

\begin{parts}
  \item Late blight of potato
  \item Loose and covered smut
  \item Transformation in bacteria
\end{parts}

\medskip


\noindent   \textbf{Question 4.} Give a brief history of discoveries in the field of microbiology focusing on the contributions made by Louis Pasteur and Robert Koch.\pts{15}

\bigskip

\begin{center}
   \textbf{SECTION -- B}
\end{center}
\medskip

\medskip


\noindent   \textbf{Question 5.} Differentiate between intergenic and intragenic interactions, and describe briefly complementary gene interaction with a suitable example.\pts{15}

\medskip


\noindent   \textbf{Question 6.} Describe briefly any two of the following:\pts{$7.5  \times 2 = 15$}

\begin{parts}
  \item Prophase I of meiosis
  \item Chromosomal theory of sex determination in plants
  \item Dominance and co-dominance
\end{parts}

\medskip


\noindent   \textbf{Question 7.} Why Mendel selected garden pea for his experiments? Briefly explain Mendel's laws of inheritance.\pts{15}

\medskip


\noindent   \textbf{Question 8.} Write short notes on any three of the following:\pts{$5  \times 3 = 15$}

\begin{parts}
  \item Spindle apparatus
  \item Crossing over and linkage maps
  \item Nuclear pore complex
  \item Chromatin
\end{parts}

\bigskip

\begin{center}
   \textbf{SECTION -- C}
\end{center}
\medskip

\medskip


\noindent   \textbf{Question 9.} Give very brief answers, not more than one sentence each, of the following:\pts{$10  \times 1 = 10$}

\begin{parts}
  \item Where does lambda DNA integrate on the chromosome of the host bacterium?
  \item Wilt of Arhar is caused by \makebox[3cm]{\hrulefill}
  \item What is the importance of denitrification in nature?
  \item What is ammonification?
  \item Replication of the bacterial chromosome.
  \item What does Hfr mean?
  \item Differentiate diakinesis from Metaphase I of meiosis.
  \item Define synthetic theory of evolution.
  \item Differentiate phenotype from genotype.
  \item What are lethal genes?
\end{parts}

\vfill
\rule{\linewidth}{0.4pt}

\begin{center}\small
\href{https://bhu-exam.vercel.app/}{Click Here}\\[4pt]\href{https://me-aryan.vercel.app/}{Click Here}\\[4pt]\href{https://quashan-me.vercel.app/}{Click Here}
\end{center}

\end{document}
